\documentclass[11pt,a4paper]{article}

\usepackage[utf8]{inputenc}
\usepackage[T1]{fontenc}
\usepackage{lmodern}
\usepackage{microtype}
\usepackage{geometry}
\geometry{top=2.0cm,bottom=2.0cm,left=2.2cm,right=2.2cm,headheight=14pt}
\usepackage{amsmath,amssymb}
\usepackage{graphicx}
\usepackage{booktabs,tabularx,array,multirow}
\usepackage{xcolor}
\usepackage{fancyhdr}
\usepackage{lastpage}
\usepackage{caption}
\usepackage{float}
\usepackage{hyperref}

\definecolor{petrol}{HTML}{1F4E5F}

\hypersetup{
  colorlinks=true,
  linkcolor=petrol,
  citecolor=petrol,
  urlcolor=petrol,
  pdftitle={Pressão de saturação e calor latente},
  pdfauthor={Rodrigo Malato Navegantes}
}

\pagestyle{fancy}
\fancyhf{}
\lhead{EMA-103 -- Laboratório de Térmica}
\rhead{Prática 2}
\cfoot{\thepage\ de \pageref{LastPage}}
\renewcommand{\headrulewidth}{0.4pt}

\setlength{\parindent}{1.2em}
\setlength{\parskip}{0.15em}
\renewcommand{\arraystretch}{1.12}
\captionsetup{font=small,labelfont=bf}
\renewcommand{\tablename}{Tabela}
\renewcommand{\figurename}{Figura}

\makeatletter
\renewcommand\section{\@startsection{section}{1}{\z@}%
  {-2.4ex \@plus -1ex \@minus -.2ex}{1.0ex \@plus .2ex}%
  {\normalfont\large\bfseries\color{petrol}}}
\renewcommand\subsection{\@startsection{subsection}{2}{\z@}%
  {-1.8ex \@plus -1ex \@minus -.2ex}{0.6ex \@plus .2ex}%
  {\normalfont\normalsize\bfseries\color{petrol}}}
\makeatother

\newcommand{\Psat}{P_{\mathrm{sat}}}
\newcommand{\Tsat}{T_{\mathrm{sat}}}
\newcommand{\hlv}{h_{lv}}
\newcommand{\Qdot}{\dot Q}
\newcommand{\mdot}{\dot m}

% Os valores abaixo são gerados por calculos_pratica2.py.
\input{resultados.tex}

\begin{document}

\begin{center}
  {\small UNIVERSIDADE FEDERAL DE MINAS GERAIS\\
  Escola de Engenharia -- Engenharia Aeroespacial\\
  EMA-103 -- Laboratório de Térmica}

  \vspace{1.25cm}
  {\LARGE\bfseries\color{petrol} Pressão de saturação e calor latente}

  \vspace{0.95cm}
  \begin{tabular}{>{\bfseries}r l}
    Aluno(a): & Rodrigo Malato Navegantes\\
    Matrícula: & 2023038175\\
    Turma: & EA1\\
    Professor(a): & Luiz Machado\\
    Data da prática: & \DataPratica\\
  \end{tabular}
\end{center}

\vspace{0.55cm}\hrule\vspace{0.65cm}

\section{Objetivo}

Determinar experimentalmente a pressão e a temperatura de saturação da água,
$\Psat$ e $\Tsat$, e o calor latente de evaporação, $\hlv$, em uma panela de
pressão tampada sem a válvula. Determinar também $\Psat$ e $\Tsat$ com a
válvula instalada e verificar a consistência dos resultados pela comparação
com valores termodinâmicos de referência, considerando as incertezas das
grandezas medidas e calculadas.

\section{Fundamentação teórica}

Durante a mudança de fase de uma substância pura em equilíbrio, a pressão e a
temperatura permanecem relacionadas pela curva de saturação. A pressão absoluta
no interior da panela é obtida pela soma entre a pressão atmosférica e a pressão
manométrica:
\begin{equation}
  P_{\mathrm{sat,exp}}=P_{\mathrm{atm}}+P_{\mathrm{man}}.
  \label{eq:psat}
\end{equation}

Com a válvula instalada, uma segunda estimativa da pressão manométrica decorre
do equilíbrio entre o peso da válvula e a força exercida pelo vapor sobre a
área do pino:
\begin{equation}
  P_{\mathrm{valv}}=\frac{m_{\mathrm{valv}}g}
  {\pi d_{\mathrm{pino}}^2/4},
  \qquad
  P_{\mathrm{sat,valv}}=P_{\mathrm{atm}}+P_{\mathrm{valv}}.
  \label{eq:valvula}
\end{equation}

Na condição sem válvula, a redução do nível da água permite calcular a vazão
mássica evaporada. Adotando a densidade $\rho_l$ da água líquida saturada na
temperatura experimental,
\begin{equation}
  \mdot_{\mathrm{evap}}=\rho_l\left(\frac{\pi d_{\mathrm{int}}^2}{4}\right)
  \frac{\Delta H}{\Delta t}.
  \label{eq:mdot}
\end{equation}

A potência fornecida ao resistor é $\Qdot_{\mathrm{el}}=V_{AB}I$. Parte dessa
potência é transferida ao ambiente pelas paredes laterais e pela tampa. Para
uma superfície lateral vertical e uma tampa horizontal aquecida, o roteiro
emprega as aproximações~\cite{roteiro}
\begin{align}
  S_v&=\pi d_{\mathrm{ext}}H,
  &S_h&=\frac{\pi d_{\mathrm{ext}}^2}{4},
  &L_c&=\frac{d_{\mathrm{ext}}}{4},
  \label{eq:areas}\\
  h_{\mathrm{conv},v}&=1{,}42
  \left(\frac{T_s-T_{\mathrm{amb}}}{H}\right)^{1/4},
  &h_{\mathrm{conv},h}&=1{,}32
  \left(\frac{T_s-T_{\mathrm{amb}}}{L_c}\right)^{1/4}.
  \label{eq:hconv}
\end{align}
Admitiu-se $T_s\simeq\Tsat$. O coeficiente linearizado de radiação é
\begin{equation}
  h_{\mathrm{rad}}=\varepsilon\sigma(T_s+T_{\mathrm{amb}})
  (T_s^2+T_{\mathrm{amb}}^2),
  \label{eq:hrad}
\end{equation}
com as temperaturas em kelvin. Assim,
\begin{align}
  \Qdot_{\mathrm{perdas}}={}&
  (h_{\mathrm{conv},v}+h_{\mathrm{rad}})S_v(T_s-T_{\mathrm{amb}})
  \nonumber\\[-0.2em]
  &+(h_{\mathrm{conv},h}+h_{\mathrm{rad}})S_h(T_s-T_{\mathrm{amb}}),
  \label{eq:qperdas}\\
  \Qdot_{\mathrm{evap}}={}&V_{AB}I-\Qdot_{\mathrm{perdas}},
  &h_{lv,\mathrm{exp}}=\frac{\Qdot_{\mathrm{evap}}}{\mdot_{\mathrm{evap}}}.
  \label{eq:hlv}
\end{align}

\section{Materiais e método}

Foram utilizados uma panela de pressão com resistor elétrico de imersão, um
termopar tipo K conectado a um medidor de temperatura, um manômetro, um
multímetro, um alicate amperímetro, uma balança, um cronômetro, uma escala de
aço e um paquímetro. A montagem é apresentada na Figura~\ref{fig:montagem}.

\begin{figure}[H]
  \centering
  \includegraphics[width=0.78\textwidth]{montagem.jpg}
  \caption{Montagem experimental da Prática 2. Fonte: material didático da
  disciplina~\cite{roteiro}.}
  \label{fig:montagem}
\end{figure}

Inicialmente, foram medidos os diâmetros interno e externo e a altura da
panela, bem como a massa da válvula e o diâmetro do pino. Na configuração sem
válvula, a panela foi aquecida até o estabelecimento de ebulição aproximadamente
estacionária. Registraram-se a pressão manométrica, a temperatura de saturação,
a tensão e a corrente elétrica. A redução do nível de água foi medida durante
o intervalo indicado pelo cronômetro. Em seguida, instalou-se a válvula e,
novamente em regime aproximadamente estacionário, registraram-se a pressão
manométrica e a temperatura de saturação.

A pressão atmosférica foi obtida na tabela da estação automática Belo
Horizonte--Pampulha (A521) do INMET~\cite{inmet}. Para 1 de setembro de 2026,
o registro das 12:00 UTC apresentou $25{,}9\,^{\circ}\mathrm{C}$, a temperatura
mais próxima da condição ambiente medida no laboratório
($\Tamb\,^{\circ}\mathrm{C}$), e pressão instantânea de $918{,}6$ hPa. Adotou-se,
portanto, $P_{\mathrm{atm}}=\Patm$ kPa.

\section{Resultados e discussão}

\subsection{Medições}

A Tabela~\ref{tab:medicoes} reúne os valores empregados no tratamento. As
incertezas indicadas nesta tabela são os limites fornecidos no roteiro ou na
apresentação das medições.

\begin{table}[H]
\centering
\caption{Grandezas medidas e respectivas incertezas.}
\label{tab:medicoes}
\small
\begin{tabular}{lccc}
\toprule
\textbf{Grandeza} & \textbf{Sem válvula} & \textbf{Com válvula} & \textbf{Incerteza}\\
\midrule
$V_{AB}$ (V) & \Vaberta & -- & $\pm3\%$\\
$I$ (A) & \Iaberta & -- & $\pm3\%$\\
$P_{\mathrm{man}}$ (bar) & \Pmanaberta & \Pmanfechada & $\pm0{,}05$ bar\\
$\Tsat$ ($^{\circ}$C) & \Tsataberta & \Tsatfechada & $\pm1{,}0\,^{\circ}$C\\
$\Delta H$ (mm) & \DeltaH & -- & $\pm2$ mm\\
$\Delta t$ (s) & \Deltat & -- & $\pm30$ s\\
$T_{\mathrm{amb}}$ ($^{\circ}$C) & \multicolumn{2}{c}{\Tamb} & $\pm1{,}0\,^{\circ}$C\\
$P_{\mathrm{atm}}$ (kPa) & \multicolumn{2}{c}{\Patm} & $\pm2\%$\\
$g$ (m/s$^2$) & \multicolumn{2}{c}{9,78} & $\pm0{,}01$ m/s$^2$\\
$m_{\mathrm{valv}}$ (g) & \multicolumn{2}{c}{\Mvalvula} & $\pm0{,}1$ g\\
$d_{\mathrm{pino}}$ (mm) & \multicolumn{2}{c}{\Dpino} & $\pm0{,}1$ mm\\
$d_{\mathrm{int}}$ (mm) & \multicolumn{2}{c}{\Dint} & $\pm1$ mm\\
$d_{\mathrm{ext}}$ (mm) & \multicolumn{2}{c}{\Dext} & $\pm1$ mm\\
$H$ (mm) & \multicolumn{2}{c}{\Hpanela} & $\pm5$ mm\\
$\varepsilon$ & \multicolumn{2}{c}{0,60} & não informada\\
\bottomrule
\end{tabular}
\end{table}

\subsection{Tratamento das incertezas}

Os limites da Tabela~\ref{tab:medicoes} foram tratados como
incertezas-padrão, as medições foram consideradas independentes e as
derivadas de sensibilidade foram calculadas numericamente. Para uma grandeza
indireta $y=f(x_1,\ldots,x_n)$, aplicou-se
\begin{equation}
  u_c(y)=\left[\sum_i\left(\frac{\partial y}{\partial x_i}u(x_i)\right)^2
  \right]^{1/2},
  \qquad U(y)=k\,u_c(y),\qquad k=2.
  \label{eq:incerteza}
\end{equation}
Logo, os resultados calculados a seguir são apresentados com incerteza
expandida, correspondente a uma abrangência de aproximadamente 95\% sob as
hipóteses usuais. Não foi incluída incerteza para a emissividade nem para os
próprios modelos de convecção, pois esses limites não foram especificados.
O programa utilizado no tratamento dos dados pode ser consultado e baixado em
\href{https://ema103-pratica2-codigo.navegantesmalato.chatgpt.site}
{Código Python da Prática 2}.

\subsection{Pressão e temperatura de saturação}

Na condição sem válvula, a Equação~\eqref{eq:psat} fornece diretamente a
pressão atmosférica, pois a leitura manométrica foi nula. Com a válvula, a
pressão devida ao peso foi
\begin{equation}
  P_{\mathrm{valv}}=\Pvalvula\pm\UPvalvula\ \mathrm{kPa}\quad(k=2).
  \label{eq:pvalv-res}
\end{equation}
A Tabela~\ref{tab:pressoes} compara as pressões absolutas experimentais com as
pressões correspondentes às temperaturas medidas. Os valores de referência
foram calculados pela equação de Antoine para a água~\cite{nist}; a incerteza
de $\Tsat$ também foi propagada nesses valores.

\begin{table}[H]
\centering
\caption{Pressões de saturação experimentais e de referência. Todas as
incertezas são expandidas, com $k=2$.}
\label{tab:pressoes}
\small
\begin{tabularx}{\textwidth}{>{\raggedright\arraybackslash}X c c c c}
\toprule
\textbf{Condição/método} &
\textbf{$P_{\mathrm{exp}}$ (kPa)} &
\textbf{$P(\Tsat)$ (kPa)} &
\textbf{Desvio} &
\textbf{Compatível}\\
\midrule
Sem válvula -- manômetro &
$\Psataberta\pm\UPsataberta$ &
$\Psattababerta\pm\UPsattababerta$ &
$+\DesvioPaberta\%$ & \CompPaberta\\
Com válvula -- manômetro &
$\Psatfechadaman\pm\UPsatfechadaman$ &
$\Psattabfechada\pm\UPsattabfechada$ &
$+\DesvioPfechadaman\%$ & \CompPfechadaman\\
Com válvula -- peso &
$\Psatfechadavalvula\pm\UPsatfechadavalvula$ &
$\Psattabfechada\pm\UPsattabfechada$ &
$+\DesvioPfechadavalvula\%$ & \CompPfechadavalvula\\
\bottomrule
\end{tabularx}
\end{table}

Os três pares são compatíveis porque seus intervalos expandidos se sobrepõem.
Na panela sem válvula, o desvio foi de $\DesvioPaberta\%$. Com a válvula, as
pressões absolutas determinadas pelo manômetro e pelo peso diferiram apenas
$\DiferencaManPeso$ kPa, valor muito menor que suas incertezas. Essa concordância
interna sustenta a leitura manométrica, embora os dois resultados tenham ficado
cerca de 6\% acima da pressão calculada a partir de $\Tsat$.

Uma verificação equivalente pode ser feita invertendo a relação de saturação,
como mostra a Tabela~\ref{tab:temperaturas}. A incerteza expandida das
temperaturas medidas é $2{,}0\,^{\circ}\mathrm{C}$.

\begin{table}[H]
\centering
\caption{Comparação entre temperaturas medidas e inferidas pelas pressões.}
\label{tab:temperaturas}
\small
\begin{tabular}{lcc}
\toprule
\textbf{Condição/método} & \textbf{$T$ medida ($^{\circ}$C)} &
\textbf{$T(P)$ ($^{\circ}$C)}\\
\midrule
Sem válvula -- manômetro & $\Tsataberta\pm2{,}0$ &
$\TsatdePaberta\pm\UTsatdePaberta$\\
Com válvula -- manômetro & $\Tsatfechada\pm2{,}0$ &
$\TsatdePfechadaman\pm\UTsatdePfechadaman$\\
Com válvula -- peso & $\Tsatfechada\pm2{,}0$ &
$\TsatdePfechadavalvula\pm\UTsatdePfechadavalvula$\\
\bottomrule
\end{tabular}
\end{table}

As temperaturas inferidas também são compatíveis com as medições. As pequenas
diferenças podem decorrer da oscilação da válvula e do manômetro, de perdas de
carga locais, do contato e posicionamento do termopar, de vazamentos pela
vedação e da aproximação de equilíbrio termodinâmico durante a fervura.

\subsection{Balanço de energia e calor latente}

Para a panela sem válvula, foram obtidos $\rho_l=\Rhol\ \mathrm{kg/m^3}$,
$S_v=\AreaV\ \mathrm{m^2}$ e $S_h=\AreaH\ \mathrm{m^2}$. Os coeficientes
estimados foram $h_{\mathrm{conv},v}=\HconvV$,
$h_{\mathrm{conv},h}=\HconvH$ e $h_{\mathrm{rad}}=\Hrad$, todos em
$\mathrm{W/(m^2K)}$. A Tabela~\ref{tab:energia} mostra as etapas seguintes do
balanço.

\begin{table}[H]
\centering
\caption{Balanço de energia e determinação do calor latente. As incertezas
calculadas são expandidas, com $k=2$.}
\label{tab:energia}
\small
\begin{tabular}{lcc}
\toprule
\textbf{Grandeza} & \textbf{Resultado} & \textbf{Unidade}\\
\midrule
Potência elétrica, $\Qdot_{\mathrm{el}}$ & $\Qeletrica\pm\UQeletrica$ & W\\
Perda pelas paredes, $\Qdot_{\mathrm{perda},v}$ & $\Qperdav\pm\UQperdav$ & W\\
Perda pela tampa, $\Qdot_{\mathrm{perda},h}$ & $\Qperdah\pm\UQperdah$ & W\\
Perda total, $\Qdot_{\mathrm{perdas}}$ & $\Qperdatotal\pm\UQperdatotal$ & W\\
Potência para evaporação, $\Qdot_{\mathrm{evap}}$ & $\Qevap\pm\UQevap$ & W\\
Massa evaporada, $m_{\mathrm{evap}}$ & $\MassaEvap$ & kg\\
Vazão evaporada, $\mdot_{\mathrm{evap}}$ & $\Mdot\pm\UMdot$ & g/s\\
Calor latente experimental, $h_{lv,\mathrm{exp}}$ &
$\Hlvexp\pm\UHlvexp$ & MJ/kg\\
Calor latente de referência, $h_{lv,\mathrm{ref}}$ &
$\Hlvtab\pm\UHlvtab$ & MJ/kg\\
\bottomrule
\end{tabular}
\end{table}

As perdas estimadas representaram $\FracPerdas\%$ da potência elétrica; assim,
a maior parte da energia fornecida foi associada à evaporação. O calor latente
experimental ficou $\DesvioHlv\%$ abaixo da referência calculada na mesma
temperatura pela correlação de Watson~\cite{watson}. Ainda assim, os intervalos
$[1{,}792;2{,}362]$ MJ/kg e $[2{,}263;2{,}275]$ MJ/kg se sobrepõem, de modo que
os valores são compatíveis no nível de abrangência adotado.

A incerteza de $\Delta H$ é relevante para a vazão evaporada, enquanto as
incertezas de tensão e corrente dominam a potência elétrica. Na estimativa pelo
peso da válvula, a dependência quadrática em $d_{\mathrm{pino}}$ torna a
incerteza desse diâmetro especialmente importante. Além desses efeitos
quantificados, podem ter contribuído para o desvio do calor latente: a
aproximação $T_s\simeq\Tsat$; o caráter empírico das correlações de convecção;
a ausência de uma incerteza atribuída à emissividade; a condução pelos apoios e
cabos; a energia armazenada caso o regime não fosse perfeitamente estacionário;
e a dificuldade de medir com precisão a variação do nível durante a ebulição.

\clearpage
\section{Conclusão}

Na panela sem válvula, foram obtidos
$\Psat=\Psataberta\pm\UPsataberta\ \mathrm{kPa}$ e
$\Tsat=\Tsataberta\pm2{,}0\,^{\circ}\mathrm{C}$. Na condição com válvula, a
pressão absoluta resultou em
$\Psatfechadaman\pm\UPsatfechadaman\ \mathrm{kPa}$ pelo manômetro e
$\Psatfechadavalvula\pm\UPsatfechadavalvula\ \mathrm{kPa}$ pelo equilíbrio do
peso. A diferença de somente $\DiferencaManPeso$ kPa entre os valores centrais
mostrou boa consistência entre os dois métodos.

Após descontar $\Qperdatotal\pm\UQperdatotal$ W de perdas térmicas, determinou-se
$h_{lv,\mathrm{exp}}=\Hlvexp\pm\UHlvexp\ \mathrm{MJ/kg}$, com incerteza expandida
$k=2$. O desvio de $\DesvioHlv\%$ em relação ao valor de referência
$\Hlvtab\pm\UHlvtab\ \mathrm{MJ/kg}$ não impediu a sobreposição das faixas de
incerteza. As pressões e temperaturas também apresentaram compatibilidade com a
relação de saturação da água. Portanto, dentro das incertezas e hipóteses do
modelo, os resultados foram coerentes e os objetivos da prática foram
alcançados.

\begin{thebibliography}{9}
\bibitem{roteiro}
UFMG. \textit{EMA-103: Laboratório de Térmica -- Aulas 4 e 5: pressão de
saturação e calor latente}. Material didático da disciplina, 2026.

\bibitem{inmet}
INSTITUTO NACIONAL DE METEOROLOGIA (INMET). \textit{Tabela de dados da estação
automática Belo Horizonte--Pampulha (A521)}. Disponível em:
\url{https://tempo.inmet.gov.br/TabelaEstacoes/A521}. Acesso em: 6 set. 2026.

\bibitem{incropera}
BERGMAN, T. L.; LAVINE, A. S.; INCROPERA, F. P.; DEWITT, D. P.
\textit{Fundamentals of Heat and Mass Transfer}. 6. ed. Hoboken: Wiley, 2011.

\bibitem{cengel}
ÇENGEL, Y. A.; BOLES, M. A. \textit{Thermodynamics: An Engineering Approach}.
8. ed. New York: McGraw-Hill, 2015.

\bibitem{nist}
NATIONAL INSTITUTE OF STANDARDS AND TECHNOLOGY (NIST). \textit{NIST Chemistry
WebBook: Water}. Disponível em:
\url{https://webbook.nist.gov/cgi/cbook.cgi?ID=C7732185\&Mask=1F}.
Acesso em: 6 set. 2026.

\bibitem{watson}
WATSON, K. M. Thermodynamics of the liquid state. \textit{Industrial \&
Engineering Chemistry}, v. 35, n. 4, p. 398--406, 1943.
\url{https://doi.org/10.1021/ie50400a004}.
\end{thebibliography}

\end{document}
